\documentclass[14pt]{extarticle}
\usepackage{fontspec}
\usepackage{polyglossia}
\usepackage{amsmath}
\usepackage{amssymb}
\usepackage{xcolor}
\usepackage{fancyhdr}
\usepackage{graphicx}
\usepackage{listings}
\usepackage[most]{tcolorbox}
% \usepackage{setspace}
\usepackage{pifont}
\usepackage{enumitem}
\usepackage{titlesec}
\usepackage[bottom]{footmisc}
\usepackage{titling}
\usepackage{minted}
\usepackage{etoolbox}
\usepackage{array}
\usepackage{longtable} % page-breakable tables (must load before polyglossia/bidi RTL setup)
\usepackage{extsizes}
\usepackage{float}
\usepackage{booktabs}
\usepackage{hyperref}
\usepackage[onehalfspacing]{setspace}
\newif\ifwhiteblack
\ifdefined\setwhiteblack
  \setwhiteblack
\fi

\ifwhiteblack
  \usepackage[a4paper, landscape, top=1.25cm, bottom=1.25cm, left=1.25cm, right=1.25cm, headheight=1.5cm, headsep=0.5cm]{geometry}
\else
  \usepackage[a4paper, top=2.7cm, bottom=2.7cm, left=2.5cm, right=2.5cm, headheight=1.5cm, headsep=0.5cm]{geometry}
\fi
\usepackage{tikz}
\usepackage{datetime2}

\usetikzlibrary{automata, positioning, arrows.meta, arrows}
\providecommand{\LTcaptype}{table}

\newfontfamily\emoji{Segoe UI Emoji}

\pagestyle{fancy}

\setmainlanguage[numerals=western]{arabic}
\setotherlanguage{english}
\newfontfamily\arabicfont[Script=Arabic]{Calibri}
\newfontfamily\arabicfonttt[Script=Arabic]{Courier New}



\ifwhiteblack
\lstset{
  language=[Sharp]C,
  numbers=left,
  stepnumber=1,
  numbersep=8pt,
  frame=single,
  basicstyle=\ttfamily\small,
  keywordstyle=\bfseries,
  stringstyle=\ttfamily,
  commentstyle=\itshape
}
\else
\lstset{
  language=[Sharp]C,
  numbers=left,
  stepnumber=1,
  numbersep=8pt,
  frame=single,
  basicstyle=\ttfamily\small,
  keywordstyle=\color{blue},
  stringstyle=\color{red},
  commentstyle=\color{green!50!black}
}
\fi

\newif\ifdetailed
\ifdefined\setdetailed
  \setdetailed
\fi

\newif\ifwithsols
\ifdefined\setwithsols
  \setwithsols
\fi

\ifwhiteblack
  \AtBeginDocument{
    \LARGE
    \usemintedstyle{bw}
  }
\fi

% Default minted look: compact code that wraps long lines.
% Not applied in black-and-white landscape mode, which already uses \LARGE.
\ifwhiteblack\else
  \setminted{fontsize=\small, breaklines=true}
\fi

% unified tcolorboxes for articles
\ifwhiteblack
\tcbset{colback=white, colframe=black, fonttitle=\bfseries, boxrule=1pt}
\newtcolorbox{boxDef}[1][]{title={{\emoji📘} تعريف\ifx\\#1\\\else ~#1\fi :}}
\newtcolorbox{boxExercise}[1][]{title={{\emoji🧩} تمرين\ifx\\#1\\\else ~#1\fi :}}
\newtcolorbox{boxExample}[1][]{title={{\emoji📝} مثال\ifx\\#1\\\else ~#1\fi :}}
\newtcolorbox{boxNote}[1][]{title={{\emoji✨} ملاحظة\ifx\\#1\\\else ~#1\fi :}}
\newtcolorbox{boxAttention}[1][]{title={{\emoji🔔} تنبيه\ifx\\#1\\\else ~#1\fi :}}
\newtcolorbox{boxWarning}[1][]{title={{\emoji⚡} ملاحظة هامة\ifx\\#1\\\else ~#1\fi :}}
\newtcolorbox{boxSolution}[1][]{title={{\emoji✅} حل\ifx\\#1\\\else ~#1\fi :}}
\newtcolorbox{boxSymbol}[1][]{title={{\emoji🔣} رمز\ifx\\#1\\\else ~#1\fi :}}
\newtcolorbox{boxHint}[1][]{title={{\emoji💡} تلميح\ifx\\#1\\\else ~#1\fi :}}
\else
\tcbset{colback=white, colframe=black, fonttitle=\bfseries, boxrule=0.8pt}
\newtcolorbox{boxDef}[1][]{colback=blue!5!white,colframe=blue!75!black,
  title={{\emoji📘} تعريف\ifx\\#1\\\else ~#1\fi :}}
\newtcolorbox{boxExercise}[1][]{colback=cyan!5!white,colframe=cyan!70!black,
  title={{\emoji🧩} تمرين\ifx\\#1\\\else ~#1\fi :}}
\newtcolorbox{boxExample}[1][]{colback=yellow!5!white,colframe=orange!90!black,
  title={{\emoji📝} مثال\ifx\\#1\\\else ~#1\fi :}}
\newtcolorbox{boxNote}[1][]{colback=gray!10!white,colframe=black,
  title={{\emoji✨} ملاحظة\ifx\\#1\\\else ~#1\fi :}}
\newtcolorbox{boxAttention}[1][]{colback=magenta!10!white,colframe=magenta!80!black,
  title={{\emoji🔔} تنبيه\ifx\\#1\\\else ~#1\fi :}}
\newtcolorbox{boxWarning}[1][]{colback=red!5!white,colframe=red!75!black,
  title={{\emoji⚡} ملاحظة هامة\ifx\\#1\\\else ~#1\fi :}}
\newtcolorbox{boxSolution}[1][]{colback=green!5!white,colframe=green!60!black,
  title={{\emoji✅} حل\ifx\\#1\\\else ~#1\fi :}}
\newtcolorbox{boxSymbol}[1][]{colback=purple!5!white,colframe=purple!70!black,
  title={{\emoji🔣} رمز\ifx\\#1\\\else ~#1\fi :}}
\newtcolorbox{boxHint}[1][]{colback=teal!5!white,colframe=teal!60!black,
  title={{\emoji💡} تلميح\ifx\\#1\\\else ~#1\fi :}}
\fi


\ifwhiteblack
\tcbset{simplecode/.style={ colback=white, colframe=black, boxrule=1pt, arc=2pt, left=4pt,right=4pt,top=4pt,bottom=4pt}}
\else
\tcbset{simplecode/.style={ colback=gray!5, colframe=black!50, boxrule=0.4pt, arc=2pt, left=4pt,right=4pt,top=4pt,bottom=4pt}}
\fi
% Compact code boxes: tighter line spacing (the preamble uses onehalfspacing) and small vertical gaps.
% Must come AFTER the simplecode definition above.
\tcbset{simplecode/.append style={fontupper=\setstretch{0.95}, top=4pt, bottom=4pt, before skip=6pt, after skip=6pt}}
\newenvironment{boxCode}{\begin{tcolorbox}[simplecode]}{\end{tcolorbox}}

% Code cell used inside tables (keeps code left-to-right inside RTL tables).
% Use inside a fixed-width p{...} column: \begin{pycell}\begin{minted}...\end{minted}\end{pycell}
\newtcolorbox{pycell}{frame empty,colback=white,boxsep=0pt,left=2pt,right=2pt,top=2pt,bottom=2pt,width=\linewidth}

\newcolumntype{C}[1]{>{\centering\arraybackslash}p{#1}}

\makeatletter
\DTMnewdatestyle{dotted}{%
  \renewcommand{\DTMdisplaydate}[4]{\two@digits{##3}.\two@digits{##2}.##1}%
  \renewcommand{\DTMDisplaydate}{\DTMdisplaydate}%
}
\makeatother
\DTMsetdatestyle{dotted}

% redefine spaces after titles
\makeatletter
\renewcommand{\@maketitle}{%
  \begin{center}
    {\huge \bfseries \@title \par}%
    \vskip 0.2em % space between title and author
    {\large \@author \par}%
    % \vskip 0.2em % space between author and date
    % {\normalsize \@date \par}%
    {\small \hfill \begin{english}\DTMtoday\end{english} \par}%
  \end{center}
}
\makeatother

\fancyhf{} % clear default
\fancypagestyle{plain}{
  \fancyhf{}
  \fancyhead[L]{مدرسة التسامح الشاملة}
  % \fancyhead[L]{\includegraphics[height=1cm]{../../../images/logoTasamoh.png}}
  \fancyhead[R]{الأستاذ محمود اغبارية}
  \fancyfoot[C]{\thepage}
}

\fancyhead[L]{مدرسة التسامح الشاملة}
\fancyhead[R]{الأستاذ محمود اغبارية}
\fancyfoot[C]{\thepage}

\setcounter{tocdepth}{3} % only section subsection and subsubsection in TOC

\newcommand{\importpdfpage}[2]{%
    \thispagestyle{empty}
    \newgeometry{left=0cm,right=0cm,top=0cm,bottom=0cm}
    \includegraphics[page=#2,width=0.95\paperwidth,height=0.95\paperheight,keepaspectratio]{#1}
    \restoregeometry
}

\newcommand{\insertFullImg}[1]{%
    \noindent
    \makebox[\textwidth][c]{
        \includegraphics[width=0.9\paperwidth,keepaspectratio]{#1}%
    }%
}%


\ifwithsols
\title{حل ورقة تمرن تمهيدية للصف العاشر $10$ - مدخل متدرّج إلى حلقة \textenglish{while}}
\else
\title{ورقة تمرن تمهيدية للصف العاشر $10$ - مدخل متدرّج إلى حلقة \textenglish{while}}
\fi

\begin{document}

\maketitle
\thispagestyle{fancy}
% ==============================================================
% القسم الأول: تتبع وفهم عمل حلقة while
% ==============================================================
\section*{القسم الأوّل: تتبّع وفهم عمل حلقة \textenglish{while}}

\begin{enumerate}[itemsep=2em]

\item تَتبَّع مقطع البرنامج التالي وأجب عن الأسئلة التي تليه:
\begin{boxCode}
\begin{english}
\begin{minted}{csharp}
int x = 2;
while (x <= 10)
{
    Console.WriteLine(x);
    x = x + 3;
}
Console.WriteLine("End: " + x);
\end{minted}
\end{english}
\end{boxCode}
\begin{enumerate}
    \item ما هي الأعداد التي يطبعها البرنامج داخل الحلقة؟
    \item ما هي قيمة \textenglish{x} النهائية التي جعلت شرط الحلقة يصبح \textenglish{false} وتتوقف الحلقة؟
\end{enumerate}
\ifwithsols
\begin{boxSolution}
\begin{enumerate}
    \item يطبع داخل الحلقة الأعداد: $2$، ثم $5$، ثم $8$.
    \item قيمة \textenglish{x} عند التوقف هي $11$ (لأنّ $11 \le 10$ غير متحقّق)، فيطبع في النهاية: \textenglish{End: 11}.
\end{enumerate}
\end{boxSolution}
\fi
\filbreak

\item تَتبَّع مقطع البرنامج التالي:
\begin{boxCode}
\begin{english}
\begin{minted}{csharp}
int n = int.Parse(Console.ReadLine());
while (n != 0)
{
    Console.WriteLine(n * 2);
    n = int.Parse(Console.ReadLine());
}
Console.WriteLine("Finished");
\end{minted}
\end{english}
\end{boxCode}
\begin{enumerate}
    \item إذا أدخل المستخدم الأعداد: $4$ ثم $7$ ثم $-3$ ثم $0$، ماذا يطبع البرنامج؟
    \item إذا كان أوّل عدد يُدخله المستخدم هو $0$، ماذا يطبع البرنامج؟
    \item لماذا نحتاج السطر \textenglish{n = int.Parse(Console.ReadLine());} في نهاية جسم الحلقة؟ ماذا سيحدث لو حذفناه؟
\end{enumerate}
\ifwithsols
\begin{boxSolution}
\begin{enumerate}
    \item يطبع: $8$، ثم $14$، ثم $-6$، ثم \textenglish{Finished} (لاحظ أنّ الـ $0$ لا يُضرب بـ $2$ لأنّ الحلقة تتوقف عنده).
    \item يطبع \textenglish{Finished} فقط، لأنّ الشرط \textenglish{n != 0} يكون \textenglish{false} منذ البداية فلا يدخل إلى جسم الحلقة أبدًا.
    \item نحتاجه لقراءة العدد التالي في كل دورة؛ لو حذفناه ستبقى قيمة \textenglish{n} الأولى ثابتة إلى الأبد ويدخل البرنامج في حلقة لا نهائية.
\end{enumerate}
\end{boxSolution}
\fi
\filbreak

\item تَتبَّع مقطع البرنامج التالي:
\begin{boxCode}
\begin{english}
\begin{minted}{csharp}
int sum = 0;
int num = int.Parse(Console.ReadLine());
while (num != -1)
{
    sum = sum + num;
    num = int.Parse(Console.ReadLine());
}
Console.WriteLine("Sum = " + sum);
\end{minted}
\end{english}
\end{boxCode}
إذا أدخل المستخدم الأعداد: $10$ ثم $5$ ثم $20$ ثم $-1$:
\begin{enumerate}
    \item أكمل جدول المتابعة التالي:
    \begin{center}
    \begin{english}
    \renewcommand{\arraystretch}{1.4}
    \begin{tabular}{|C{2.5cm}|C{3.5cm}|C{2.5cm}|C{3.5cm}|}
    \hline
    \textbf{num} & \textbf{num != -1} & \textbf{sum} & \textarabic{\textbf{المخرجات}} \\
    \hline
     & & $0$ & \\ \hline
    $10$ & \textenglish{true} & & \\ \hline
    $5$ & & & \\ \hline
    $20$ & & & \\ \hline
    $-1$ & & & \\ \hline
    \end{tabular}
    \end{english}
    \end{center}
    \item هل جُمِع العدد $-1$ مع المتغير \textenglish{sum}؟ لماذا؟
\end{enumerate}
\ifwithsols
\begin{boxSolution}
\begin{enumerate}
    \item .
    \begin{center}
    \begin{english}
    \renewcommand{\arraystretch}{1.4}
    \begin{tabular}{|C{2.5cm}|C{3.5cm}|C{2.5cm}|C{3.5cm}|}
    \hline
    \textbf{num} & \textbf{num != -1} & \textbf{sum} & \textarabic{\textbf{المخرجات}} \\
    \hline
     & & $0$ & \\ \hline
    $10$ & \textenglish{true} & $10$ & \\ \hline
    $5$ & \textenglish{true} & $15$ & \\ \hline
    $20$ & \textenglish{true} & $35$ & \\ \hline
    $-1$ & \textenglish{false} & $35$ & \textenglish{Sum = 35} \\ \hline
    \end{tabular}
    \end{english}
    \end{center}
    \item كلا، لم يُجمع العدد $-1$ لأنّنا قرأناه في نهاية الدورة الثالثة، وعندما فُحص الشرط \textenglish{num != -1} أعاد \textenglish{false} فخرجت الحلقة فورًا قبل تنفيذ \textenglish{sum = sum + num}.
\end{enumerate}
\end{boxSolution}
\fi
\filbreak

\end{enumerate}

\clearpage
% ==============================================================
% القسم الثاني: كتابة شرط الاستمرار وإكمال الكود الناقص
% ==============================================================
\section*{القسم الثاني: كتابة شرط الاستمرار وإكمال الكود الناقص}

\begin{enumerate}[resume, itemsep=2em]

\item \textbf{تدريب على كتابة شرط \textenglish{while}:} \\
تذكّر أنّنا نكتب داخل \textenglish{while (...)} \textbf{شرط الاستمرار} (متى تستمرّ الحلقة؟)، وهو \textbf{عكس شرط التوقف}. \\
لكل حالة في الجدول التالي، اكتب الشرط المناسب داخل \textenglish{while (...)}:

{\centering
\renewcommand{\arraystretch}{1.6}
\begin{tabular}{|>{\raggedleft\arraybackslash}p{0.62\textwidth}|>{\centering\arraybackslash}p{0.32\textwidth}|}
\hline
\textbf{المطلوب في السؤال (متى تتوقف الحلقة؟)} & \textbf{شرط الاستمرار في \textenglish{while}} \\
\hline
$1$. تتوقف الحلقة عندما يُدخل المستخدم العدد $0$ في المتغير \textenglish{n}. & \textenglish{while (n != 0)} \\
\hline
$2$. تتوقف الحلقة عندما يُدخل المستخدم العدد $-1$ في المتغير \textenglish{x}. & \textenglish{while (............)} \\
\hline
$3$. تتوقف الحلقة عندما يُدخل المستخدم عددًا \textbf{سالبًا} في المتغير \textenglish{num} (أي تستمر طالما العدد غير سالب). & \textenglish{while (............)} \\
\hline
$4$. تتوقف الحلقة عندما يُدخل المستخدم كلمة المرور الصحيحة $1234$ في المتغير \textenglish{pass}. & \textenglish{while (............)} \\
\hline
$5$. تتوقف الحلقة عندما يصبح المجموع \textenglish{sum} أكبر تمامًا من $100$. & \textenglish{while (............)} \\
\hline
$6$. تتوقف الحلقة عندما يُدخل المستخدم في \textenglish{n} عددًا يقبل القسمة على $7$ بدون باقٍ. & \textenglish{while (............)} \\
\hline
\end{tabular}
\par}
\ifwithsols
\begin{boxSolution}
\begin{enumerate}[nosep]
    \item \textenglish{while (n != 0)}
    \item \textenglish{while (x != -1)}
    \item \textenglish{while (num >= 0)}
    \item \textenglish{while (pass != 1234)}
    \item \textenglish{while (sum <= 100)}
    \item \textenglish{while (n \% 7 != 0)}
\end{enumerate}
\end{boxSolution}
\fi
\filbreak

\item \textbf{المستوى $1$ — أكمل شرط الحلقة فقط:}
\begin{enumerate}
    \item البرنامج التالي يستقبل أعدادًا صحيحة ويطبع كل عدد يدخله المستخدم، ويتوقف عندما يُدخل المستخدم العدد $0$. أكمل الشرط الناقص:
\begin{boxCode}
\begin{english}
\begin{minted}{csharp}
int n = int.Parse(Console.ReadLine());
while ( ____________ )
{
    Console.WriteLine(n);
    n = int.Parse(Console.ReadLine());
}
\end{minted}
\end{english}
\end{boxCode}

    \item البرنامج التالي يطلب من المستخدم إدخال عدد موجب (أكبر من $0$). طالما أدخل المستخدم عددًا أصغر أو يساوي $0$، يطبع البرنامج رسالة خطأ ويطلب عددًا جديدًا:
\begin{boxCode}
\begin{english}
\begin{minted}{csharp}
int x = int.Parse(Console.ReadLine());
while ( ____________ )
{
    Console.WriteLine("Error! Enter a positive number:");
    x = int.Parse(Console.ReadLine());
}
Console.WriteLine("Thank you! You entered: " + x);
\end{minted}
\end{english}
\end{boxCode}
\end{enumerate}
\ifwithsols
\begin{boxSolution}
\begin{enumerate}[nosep]
    \item الشرط هو: \textenglish{n != 0}
    \item الشرط هو: \textenglish{x <= 0}
\end{enumerate}
\end{boxSolution}
\fi
\filbreak

\item \textbf{المستوى $2$ — أكمل الشرط وسطْر التحديث داخل الحلقة:}
\begin{enumerate}
    \item البرنامج التالي يطلب من المستخدم إدخال عدد يقبل القسمة على $7$ بدون باقٍ، ويكرّر الطلب حتى يُدخل المستخدم عددًا يقبل القسمة على $7$ ثم يطبعه:
\begin{boxCode}
\begin{english}
\begin{minted}{csharp}
int n = int.Parse(Console.ReadLine());
while ( ____________ )
{
    n = ____________________________;
}
Console.WriteLine( ________________ );
\end{minted}
\end{english}
\end{boxCode}

    \item البرنامج التالي يستقبل كلمة المرور (عدد صحيح)، وطالما كانت غير مطابقة لكلمة المرور الصحيحة $8264$ يطلب إدخالها من جديد، وفي النهاية يطبع \textenglish{"Welcome"}:
\begin{boxCode}
\begin{english}
\begin{minted}{csharp}
int pass = int.Parse(Console.ReadLine());
while ( ____________ )
{
    Console.WriteLine("Wrong password, try again:");
    ____________________________________;
}
Console.WriteLine("Welcome");
\end{minted}
\end{english}
\end{boxCode}
\end{enumerate}
\ifwithsols
\begin{boxSolution}
\begin{english}
\begin{minted}{csharp}
// a
int n = int.Parse(Console.ReadLine());
while (n % 7 != 0)
{
    n = int.Parse(Console.ReadLine());
}
Console.WriteLine(n);

// b
int pass = int.Parse(Console.ReadLine());
while (pass != 8264)
{
    Console.WriteLine("Wrong password, try again:");
    pass = int.Parse(Console.ReadLine());
}
Console.WriteLine("Welcome");
\end{minted}
\end{english}
\end{boxSolution}
\fi
\filbreak

\item \textbf{المستوى $3$ — أكمل حلقات تحتوي على عدّاد (\textenglish{count}) أو مجموع (\textenglish{sum}):}
\begin{enumerate}
    \item البرنامج التالي يستقبل أعدادًا صحيحة حتى يُدخل المستخدم العدد $0$، وفي النهاية يطبع \textbf{عدد} الأعداد التي أُدخلت (بدون حساب الـ $0$):
\begin{boxCode}
\begin{english}
\begin{minted}{csharp}
int count = ____;
int n = int.Parse(Console.ReadLine());
while ( ____________ )
{
    count = ____________;
    n = int.Parse(Console.ReadLine());
}
Console.WriteLine( ________________ );
\end{minted}
\end{english}
\end{boxCode}

    \item البرنامج التالي يستقبل أعدادًا صحيحة ويجمعها، ويتوقف عندما يصبح المجموع أكبر تمامًا من $100$، ثم يطبع المجموع النهائي:
\begin{boxCode}
\begin{english}
\begin{minted}{csharp}
int sum = ____;
while ( ____________ )
{
    int v = int.Parse(Console.ReadLine());
    sum = ____________;
}
Console.WriteLine( ________________ );
\end{minted}
\end{english}
\end{boxCode}

    \item البرنامج التالي يستقبل أعدادًا صحيحة حتى يُدخل المستخدم $-1$، ويعدّ ويطبع في النهاية كم عددًا من بينها كان \textbf{موجبًا} (أكبر من $0$):
\begin{boxCode}
\begin{english}
\begin{minted}{csharp}
int posCount = 0;
int n = int.Parse(Console.ReadLine());
while ( ____________ )
{
    if ( ____________ )
    {
        posCount++;
    }
    n = ________________________________;
}
Console.WriteLine( ________________ );
\end{minted}
\end{english}
\end{boxCode}
\end{enumerate}
\ifwithsols
\begin{boxSolution}
\begin{english}
\begin{minted}{csharp}
// a
int count = 0;
int n = int.Parse(Console.ReadLine());
while (n != 0)
{
    count = count + 1; // or count++;
    n = int.Parse(Console.ReadLine());
}
Console.WriteLine(count);

// b
int sum = 0;
while (sum <= 100)
{
    int v = int.Parse(Console.ReadLine());
    sum = sum + v; // or sum += v;
}
Console.WriteLine(sum);

// c
int posCount = 0;
int n = int.Parse(Console.ReadLine());
while (n != -1)
{
    if (n > 0)
    {
        posCount++;
    }
    n = int.Parse(Console.ReadLine());
}
Console.WriteLine(posCount);
\end{minted}
\end{english}
\end{boxSolution}
\fi
\filbreak

\end{enumerate}


\end{document}
